% !TeX root = ../../Book4.tex
% 纲要标题：## 附录 E　作用滤过与持续 Floer 结构
% 纲要位置：几何学计划纲要.md，第 4685 行。
% 正文状态：未撰写；本文件目前是框架。

% 纲要内容开始
% 
% > 建立有限型过滤复形的代数理论，再与正文受限 Hamilton 模块连接。
% 
% * **过滤链复形**
%   * 作用值、子复形与商复形
%   * 延拓造成的过滤移位
%   * 谱数的定义和存在
% 
% * **持续模与稳定性**
%   * 有限型区间分解
%   * 交错和瓶颈距离
%   * 条形码稳定性所需条件
%   * 无限维与 Novikov 情形的修正
% 
% * **Hamilton 距离**
%   * Hofer 长度及范数
%   * 谱不变量给出的估计
%   * 具体平移和支撑模型
% 
% ---
% 
% 纲要内容结束

\chapter{作用滤过与持续 Floer 结构}
\label{g4:app:E}
\BookChapterNumber{g4:app:E}

待撰写

\section{过滤链复形}
\label{g4:sec:E01}
待撰写

\subsection{作用值、子复形与商复形}
\label{g4:subsec:E01:01}
待撰写

\subsection{延拓造成的过滤移位}
\label{g4:subsec:E01:02}
待撰写

\subsection{谱数的定义和存在}
\label{g4:subsec:E01:03}
待撰写

\section{持续模与稳定性}
\label{g4:sec:E02}
待撰写

\subsection{有限型区间分解}
\label{g4:subsec:E02:01}
待撰写

\subsection{交错和瓶颈距离}
\label{g4:subsec:E02:02}
待撰写

\subsection{条形码稳定性所需条件}
\label{g4:subsec:E02:03}
待撰写

\subsection{无限维与 Novikov 情形的修正}
\label{g4:subsec:E02:04}
待撰写

\section{Hamilton 距离}
\label{g4:sec:E03}
待撰写

\subsection{Hofer 长度及范数}
\label{g4:subsec:E03:01}
待撰写

\subsection{谱不变量给出的估计}
\label{g4:subsec:E03:02}
待撰写

\subsection{具体平移和支撑模型}
\label{g4:subsec:E03:03}
待撰写
